\subsection{Comparing complete certificates at the remaining odd roots}
\label{sec:odd-complete:comparison}

All operators and vectors in this section use the original finite
probability measure. Write $z=i\sqrt x$, $\theta=13/32$,
$\kappa=19/32$, $\tau=d-\delta_2$, and let $\nu$ have mass $d$.
We compare the same actual positive terminal majorant
\begin{equation}\label{eq:odd-complete:input}
 X\le e^\alpha[Cq+c+BQ(B)u],\qquad
 \gamma=C-1,\quad k_0=C+c-1.
\end{equation}
The coefficients $C,c$ and all moments of $Q$ retain their actual
values. In particular $\gamma\one$ is a vector of norm $|\gamma|$.

\paragraph{The finite bare-exponent packet.}
For an actual $F\le e^{R-\Gamma}$, with $R=r(B)u$, put
$D_i=\E B^iR$, $\delta=\E R+c_T\E R^2$ and
$\epsilon=\delta+D_1$, using the pointwise Taylor coefficient
$c_T$ proved for that terminal.
The positive first row and its finite continuation give
\[
 e^\Gamma T^LF\le\bar f_L+\delta\bar f_{L-1}+B^LR
                    +\sum_{i=1}^{L-1}D_i\bar f_{L-1-i}.
\]
Substituting the exact short paths into this same majorant gives
the parameters of \eqref{eq:odd-complete:input}, with $\alpha=-\Gamma$:
\begin{equation}\label{eq:odd-even:rows}
\begin{array}{c|c|c|c}
 L&C&c&Q\\\hline
 1&1&\delta&r\\
 2&1+\delta&D_1&1+zr\\
 3&1+\epsilon&-d+D_2&G_2+\delta+z^2r\\
 \ge4&1-hd&2\delta_++2(D_1)_+-hd&G_{L-1}+z^{L-1}r.
\end{array}
\end{equation}
The last row assumes $\gamma_i=\beta_i-D_{i-1}\ge hd$ also
at the mean index $i=L+1$, and retains the final two paths
$\bar f_0,\bar f_1$ from that positive convolution. The signed
means in the four windows are exactly
\begin{equation}\label{eq:odd-complete:bare-means}
 D_1,\qquad -d+D_2,\qquad -(1+\delta)d+D_3,\qquad
 -\gamma_{L+1}.
\end{equation}
Both the constant $C-1$ and these moments belong to the same $R$.

\begin{lemma}[One certificate with an explicit constant split]
\label{lem:odd-complete:certificate}
Suppose an actual root pair satisfies
$P\le\exp(BS-\zeta u-\Gamma_r)$, and its separate unit arm is $p$.
For any real $\eta$ put
\[
 V=Q(B)u+S+(\gamma-\zeta-\eta)\one.
\]
Then $H=pPX$ has the pointwise certificate
\begin{equation}\label{eq:odd-complete:certificate}
 \log H\le\alpha+k_0-\Gamma_r+BV+\eta u-3u^2/10,
\end{equation}
and its complete fourth exponent is at most
\[
 4(\alpha+k_0-\Gamma_r+\E BV)
       +(17/4)\|V\|^2+(10/3)\eta^2.
\]
For the auxiliary pair $P_0\le e^{-d-B^2u}$, the choices
$\eta=\gamma$ and $\eta=0$ give, respectively, the curved and
full-constant baselines
\begin{align}
 J_0^c&=4(\alpha+k_0+\E BQu-d)
        +(17/4)\|(Q-B)u\|^2+(10/3)\gamma^2,
       \label{eq:odd-complete:baseline-curved}\\
 J_0^f&=4(\alpha+k_0+\E BQu-d)
        +(17/4)\|(Q-B)u+\gamma\one\|^2.
       \label{eq:odd-complete:baseline-full}
\end{align}
\end{lemma}
\begin{proof}
Apply $\log y\le y-1$ to the positive bracket in
\eqref{eq:odd-complete:input}, and add the actual pair envelope and
$\log p\le-u-3u^2/10$. Since $B\one=u$, this gives
\eqref{eq:odd-complete:certificate} exactly. Apply
Lemma~\ref{lem:odd-common:fourth} to these complete vectors.
\end{proof}

\begin{lemma}[Minimum-one root transfers]
\label{lem:odd-complete:minimum-one}
Take $3\le a<b$ and the actual root $pW_a\bar f_b$.
Write $Q=A+zD$. Suppose $A\ge0$, $D\le\bar D$, $\bar D\ge1$,
$\|Q(B)u\|\le R\sqrt d$, and $\gamma\ge-\chi d$ on
$0<d\le d_*$, with $0\le\chi\le17/10$, $d_*\le1/3$.
Relative to the curved baseline, the complete exponent increments
are at most
\begin{align}
 J_3-J_0^c&\le d\Phi(\bar D,\chi,d_*),
    \label{eq:odd-complete:Phi}\\
 J_{a\ge4}-J_0^c&\le d\Psi(\bar D,R),
    \label{eq:odd-complete:Psi}
\end{align}
where
\[
\begin{aligned}
 \Phi(\bar D,\chi,d_*)&=-8+
 [12+(17/4)(2\bar D-1+\theta)]\theta
       +(10/3)d_*[\kappa^2+2\chi\kappa],\\
 \Psi(\bar D,R)&=-127/15+
       (17/4)[(2\bar D-1)\theta+3R/4+9/64].
\end{aligned}
\]
For a full-constant baseline with $\gamma\ge-\chi d$, use instead
$\Phi(\bar D,0,d_*)$ or $\Psi(\bar D,R)$, and add at most
\begin{equation}\label{eq:odd-complete:full-fees}
 (17/6)\chi d\quad(a=3),\qquad
 (17/6)\chi(1+\theta^2)d\quad(a\ge4).
\end{equation}
If $A\ge-\rho_0$ replaces $A\ge0$, replace
$2\bar D-1$ by $2\bar D-1+2\rho_0$ in $\Phi$.
\end{lemma}
\begin{proof}
The actual reversed length-three pair, including its constant loss,
is Lemma~\ref{lem:odd-common:root-three} in the form
\[
 W_3\bar f_b\le e^{BS-\tau u-(3d-2\delta_2)},\qquad
 S=-x(E+zF)u,
 \quad0\le E,F\le1.
\]
The scalar plus signed-mean increment is at most
$-8d+12\delta_2$. The nonconstant norm differences satisfy
\[
 (A-xE)^2-A^2\le x^2,
 \qquad
 x[(D-xF)^2-(D-1)^2]
   =x(1-xF)(2D-1-xF)\le(2\bar D-1)x.
\]
The curved choice of Lemma~\ref{lem:odd-complete:certificate}
changes $\gamma$ to $\gamma-\tau$, costing at most
$(10/3)(\tau^2+2\chi d\tau)$.
Only after integrating set $t=\delta_2/d$. The resulting upper
\[
 -8+[12+(17/4)(2\bar D-1+\theta)]t
       +(10/3)d_*[(1-t)^2+2\chi(1-t)]
\]
increases on $[0,\theta]$: its derivative is at least
$12+(17/4)(1+\theta)-(20/9)(1+17/10)>0$.
This proves $\Phi$ without replacing the same global moment twice.
For $A\ge-\rho_0$, its real norm difference adds $2\rho_0x$.

For $a\ge4$, the actual pair estimate \eqref{eq:L96:long} gives
\[
 W_a\bar f_b\le e^{BS-16d/5},\qquad
 S=\sum_{j=a+1}^{b}B^{j-2}u.
\]
Its contiguous finite mask gives
$\|S\|<3\sqrt d/8$, $\E BS\le d/12$, and
$-\delta_3\le\E S\le\delta_2$.
Thus the scalar plus mean improves by at least $127d/15$.
The base norm difference is
$\|Qu\|^2-\|(Q-B)u\|^2=\int(2D-1)x\,d\nu$;
the further norm increment is at most
$2R(3/8)d+9d/64$. The curvature coefficient is unchanged,
giving $\Psi$.

For the full-constant convention choose $\eta=-\tau$ at $a=3$
and $\eta=0$ at $a\ge4$. The additional exact cross term is
$2\gamma\E(S+Bu)$, while
$\E(S+Bu)\in[-d,-\tau]$ or
$[-d-\delta_3,-\tau]$, respectively. A nonnegative $\gamma$
helps; otherwise $\gamma\ge-\chi d$ and $d\le1/3$ give
\eqref{eq:odd-complete:full-fees}, with the constant's signed mean
retained inside both vectors.
\end{proof}

\begin{lemma}[Positive-root transfers, with signed overshoots]
\label{lem:odd-complete:positive}
Take the actual root $f_a h_b f_c$, $3\le a<b<c$.
Write $G_{a-1}(-z)=E_a-zD_a$ and
$S_{bc}=(r+zs)u=\sum_{j=b+1}^{c}(-1)^jB^{j-2}u$.
Put $A_*=E_a+r$, $\phi=1-D_a+s$, and
$V_*=(Q+G_{a-1}(-B))u+S_{bc}$.
For $Q=U+zV$ define
\[
 \mathcal T(Q)=\frac{17}{2}\int
 [(1+x^2)(U-1)_+ +(x^2+x^3)(V-1)_+ +x^3(-V)_+]\,d\nu.
\]
The complete exponent differences are bounded by
\begin{align}
 J_+-J_0^f&\le-11d/10+\mathcal T(Q)
                 -(17/2)\gamma\int\phi\,d\nu,
       \label{eq:odd-complete:positive-full}\\
 J_+-J_0^c&\le-11d/10+\mathcal T(Q)
                 +(17/2)\gamma\E V_*+(11/12)\gamma^2.
       \label{eq:odd-complete:positive-curved}
\end{align}
Alternatively, if $Q=Q_0+\Delta Q$ with both normalized coordinates
of $Q_0$ in $[0,1]$, replace $-11d/10+\mathcal T(Q)$ by
\begin{equation}\label{eq:odd-complete:positive-norm}
 -d+(34/3)\sqrt d\,\|\Delta Q(B)u\|.
\end{equation}
The same full-constant or curved correction still applies.
\end{lemma}
\begin{proof}
Subtract the actual exponents in
Lemma~\ref{lem:odd-common:replacement}. This gives its exact density
$j_Q$, together with the displayed two distinct constant corrections.
On $0\le U,V\le1$ the three actual finite root blocks have upper
polynomials
\[
\begin{array}{c|l}
 (a,b)&j_Q\text{ upper}\\\hline
 (3,4)&-81/20+13x^2+(17/4)x^3+(17/4)x^4,\\
 (3,b\ge5)&-81/20+13x^2+(17/2)x^3+(17/4)x^4+(17/4)x^5,\\
 a\ge4&-81/20-5x+22x^2+(17/2)x^3+(17/4)x^4+(17/4)x^5.
\end{array}
\]
For the first block $E_a=D_a=1$, $r=x^2F$, $s=xE$,
with $0\le E,F\le1$, including the finite gap one.
For $b\ge5$ the four starting residues give $|r|,|s|\le x^2$.
Convexity selects $r=x^2$, then $U=1$; the imaginary part is
affine in $V$ and convex in $D_a-s$.
For $a\ge4$, $E_a\le1-x+x^2$, $D_a\ge1-x$; decreasing $E_a$
from one saves at least $9x(1-x)$, while relaxing $-4D_a$ costs
at most $4x$. These give the three displayed polynomials.
The first two are below $-11/10$ using
$\theta^2<1/6$, $\theta^3<1/14$, $\theta^4<1/36$.
The third is convex on $[0,5/12]$; at its right endpoint its
excess over $-3/2$ is $-87307/4976640$.

All finite blocks also give
$0<A_*\le1+x^2$, $-x^2\le\phi\le x+x^2$.
The dependence of the exact density on the terminal coordinates is
\[
 \Delta j=(17/2)(A_*\Delta U+x\phi\Delta V).
\]
Clipping $Q$ into the unit rectangle therefore gives the three
signed charges in $\mathcal T$. A negative real coordinate only
helps. For the norm version the coefficient is at most
$(17/2)\sqrt{A_*^2+x\phi^2}<34/3$; use $|\phi|\le2x$ and
$(1+\theta^2)^2+4\theta^3<16/9$, then Cauchy with the same
spectral measure of mass $d$. The $L^2$ tail is used in this
integral Cauchy inequality.
\end{proof}

\begin{corollary}[Bare comparisons in the same directions]
\label{cor:odd-complete:bare}
For $n\ge3$ the actual $pW_a\bar f_b\bar f_n$, $3\le a<b$,
has fourth loss at least $29/10$.
For the positive root $f_a h_bf_c$, $3\le a<b<c$, the actual
bare input $\bar f_n$ has respective losses
\[
 419/128\ (n=1),\quad387/128\ (n=2),\quad35/4\ (n=3),
 \quad11/2\ (n\ge4).
\]
\end{corollary}
\begin{proof}
Retain $\beta_n\bar f_0$ in the positive convolution:
$\bar f_n\le q-\beta_n+BG_{n-1}(B)u$.
At the auxiliary root the actual exponents have losses
$291/128$, $259/128$, $31/4$, and at least $9/2$, respectively.
For $n\ge4$ this uses
$\beta_n+\beta_{n+1}\ge2\kappa d$ and
$|G_{n-1}-z|^2\le(1-x+x^2)^2+x^3\le1$.
Every $G_{n-1}$ is in the coordinate rectangle.
The positive transfer saves at least $d$.
For the minimum-one root, $\gamma=0$, $\bar D=1$, $R\le6/5$;
both $\Phi(1,0,1/3)$ and $\Psi(1,6/5)$ are negative.
The weaker claimed $29/10$ follows.
\end{proof}
